\section{Turning the average argument into a construction}
The random proof shows that a favorable labeling exists. We can also construct one by deciding the labels sequentially while keeping the expected final cut size from decreasing.

At a partial assignment, let $C$ be the number of crossing edges whose endpoints are both already assigned. Let $U$ be the number of edges with at least one unassigned endpoint. If all remaining labels were chosen independently and fairly, every unresolved edge would cross with probability $1/2$, whether one or both endpoints are unassigned. The expected final count is therefore
\[
F=C+\frac U2.
\]
At the beginning, $F=m/2$.

Now assign the next vertex $v$. Suppose it has $\ell$ already assigned neighbors on the left and $r$ on the right. Sort the edges at $v$ into three kinds and see how each one's contribution to $F$ changes if we place $v$ on the left.
\begin{itemize}
\item The $\ell$ edges to assigned neighbors on the left were unresolved, contributing $1/2$ each. Afterwards both endpoints are on the left, so they do not cross and contribute $0$.
\item The $r$ edges to assigned neighbors on the right also contributed $1/2$ each. Afterwards they cross and contribute $1$ each.
\item Edges from $v$ to vertices that are still unassigned remain unresolved and still contribute $1/2$ each.
\end{itemize}
Edges not touching $v$ are unaffected. So the change in $F$ is
\[
r\cdot\Big(1-\frac12\Big)+\ell\cdot\Big(0-\frac12\Big)=\frac{r-\ell}{2}.
\]
Placing $v$ on the right instead gives change $(\ell-r)/2$. These two changes are negatives of each other, so at least one of them is nonnegative. The rule is therefore simple: put $v$ on the side \emph{opposite} the larger number of its already assigned neighbors (either side if the numbers are equal). Then $F$ never decreases.

After all vertices are assigned, $U=0$ and $F$ is the actual cut size. Since $F$ never decreased from $m/2$, the final cut has at least $m/2$ edges. This is a deterministic construction with a maintained lower bound. The technique is an instance of conditional expectation, but its entire calculation is visible here; no general probability theory is needed.

\subsection*{A full trace of the maintained average}
Take a triangle with vertices $a,b,c$ and three edges. Initially no labels are assigned, so $C=0$, $U=3$, and $F=3/2$. Assign $a$ to the left. No edge has both endpoints assigned yet, so $F$ stays $3/2$.

Vertex $b$ has one assigned neighbor on the left and none on the right: $\ell=1,r=0$. Put $b$ on the right. Edge $ab$ is now a crossing edge, while $ac$ and $bc$ remain unresolved. Thus $C=1,U=2$, and $F=2$. The increase is $(\ell-r)/2=1/2$, as predicted.

Vertex $c$ has one assigned neighbor on each side, so either choice changes $F$ by zero. Put it on the left. Now $ab$ and $bc$ cross, while $ac$ does not. We have $C=2,U=0,F=2$. The final value is an actual cut size, not an expectation about unfinished choices.

This trace distinguishes three quantities. $C$ counts crossings already determined. $U$ counts edges still awaiting at least one label. $F$ estimates the final count under stipulated fair choices for the unassigned labels. The algorithm need not actually toss coins. At each step it chooses a side whose computed $F$ is at least the current value. When nothing remains unassigned, the maintained estimate becomes the actual answer.
